\chapter{Reproducibility Checklist}
\label{app:reproducibility}

This appendix records the operational material behind the observed
platform configuration in \secref{sec:cresco8-observed}.  It is
separated from the platform chapter because the commands support the
argument but are not part of it.  The hardware measurement campaigns,
with their kernels, scripts and raw records, and the platform record
described below are maintained in a separate repository, which is the
single source for both~\cite{hwbench2026}.

\section{Configuration probes}
\label{app:configuration-probes}

Every item of \tabref{tab:cresco8-manifest} is read by an ordinary user
inside an exclusive allocation, without root privileges, on the
execution node itself.  The benchmark suite's system snapshot performs
these reads at the start of every campaign and stores them as a
structured platform record next to the topology map, together with the
\texttt{lstopo} capture of \figref{fig:cresco8-lstopo}.
\tabref{tab:configuration-probes} lists the source of each item, so that
the observation can be repeated on any later date or node.

\begin{table}[htbp]
\centering
\small
\tabsetup
\begin{tabularx}{\textwidth}{P{3.0cm}Y}
\headrow
\headcell{Item} & \headcell{Command or kernel interface} \\[3pt]
\tabheadrule
\tabstripeRow
Topology
 & \texttt{lstopo}, \texttt{numactl -{}-hardware}, \texttt{lscpu},
   \texttt{lscpu -C} \\[3pt]

\tabrowsep
\tabplainRow
Memory-side cache
 & \texttt{/sys/devices/system/node/node*/\allowbreak memory\_side\_cache/\allowbreak index*/}
   \{\texttt{size}, \texttt{line\_size}, \texttt{indexing},
   \texttt{write\_policy}\} \\[3pt]

\tabrowsep
\tabstripeRow
DIMM population
 & \texttt{/sys/devices/system/edac/\allowbreak mc/mc*/dimm*/}\{\texttt{size},
   \texttt{dimm\_label}\} \\[3pt]

\tabrowsep
\tabplainRow
Threads and frequency
 & \texttt{/sys/devices/system/cpu/smt/active};
   \texttt{/sys/devices/system/cpu/cpu0/\allowbreak cpufreq/}\{\texttt{scaling\_driver},
   \texttt{scaling\_governor}\}; \texttt{/proc/cmdline} \\[3pt]

\tabrowsep
\tabstripeRow
Page policy
 & \texttt{/proc/sys/kernel/numa\_balancing},
   \texttt{/proc/sys/vm/zone\_reclaim\_mode},
   \texttt{/sys/kernel/mm/numa/demotion\_enabled},
   \texttt{/sys/kernel/mm/transparent\_hugepage/}\{\texttt{enabled},
   \texttt{defrag}\} \\[3pt]

\tabrowsep
\tabplainRow
Counters and energy
 & \texttt{/proc/sys/kernel/perf\_event\_paranoid};
   \texttt{/sys/class/powercap/intel-rapl:0/energy\_uj};
   \texttt{/dev/cpu/0/msr}; \texttt{scontrol show node} for Slurm's
   energy fields \\[3pt]

\tabrowsep
\tabstripeRow
Job cgroup
 & \texttt{cpuset.mems.effective} under the job's path in
   \texttt{/sys/fs/cgroup}, found from \texttt{/proc/self/cgroup} \\[3pt]

\tabrowsep
\tabplainRow
Network and storage
 & \texttt{ibstat}; \texttt{lsblk}; \texttt{df -hT} on \texttt{/tmp},
   \texttt{/dev/shm} and the Lustre mount \\[3pt]

\tabrowsep
\tabstripeRow
Scheduler
 & \texttt{sinfo}, \texttt{scontrol show partition},
   \texttt{scontrol show config}, \texttt{sacctmgr show qos} \\[3pt]
\tabbottomrule
\end{tabularx}
\caption{Sources of the observed configuration items.  All are readable
without root privileges; the RAPL and MSR entries are listed because
their \emph{unreadability} is the observation.}
\label{tab:configuration-probes}
\end{table}

\section{Minimum platform record}
\label{app:platform-record}

Each reported measurement carries, in addition, the record in
\tabref{tab:platform-record}.

\begin{table}[htbp]
\centering
\small
\tabsetup
\begin{tabularx}{\textwidth}{lY}
\headrow
\headcell{Category} & \headcell{Record} \\[3pt]
\tabheadrule
\tabstripeRow
Identity
 & Hostname, Slurm job ID, partition, date and node allocation. \\[3pt]

\tabrowsep
\tabplainRow
Processor
 & Exact SKU, sockets, physical cores, SMT state, microcode, governor
   and frequency sampled during the run. \\[3pt]

\tabrowsep
\tabstripeRow
Memory
 & Memory and clustering modes, NUMA distance matrix, HBM and DDR
   node IDs, huge-page state and verified page residency. \\[3pt]

\tabrowsep
\tabplainRow
I/O
 & Location of the model file and how it was loaded (staged copy,
   memory map or anonymous memory). \\[3pt]

\tabrowsep
\tabstripeRow
Software
 & OS and kernel, loaded modules, compiler, runtime commit, linked
   libraries, build flags and complete command line. \\[3pt]

\tabrowsep
\tabplainRow
Telemetry
 & Counter and energy interfaces used, with their permissions and
   measurement boundaries. \\[3pt]
\tabbottomrule
\end{tabularx}
\caption{Minimum platform record for an inference result.}
\label{tab:platform-record}
\end{table}

\section{Exposing HBM in Flat mode}
\label{app:flat-mode-setup}

Flat mode needs two configuration stages on Linux.  The firmware must
first select \emph{Flat} under \emph{Memory Hierarchy}.  After boot, the
HBM is presented as DAX devices rather than as memory, and each device
must be converted to system RAM before it appears as an allocatable
memory-only NUMA node~\cite{lp1738}:

\begin{lstlisting}[style=shell,
  caption={Converting the HBM DAX devices of a Flat-mode node to system
  RAM.  Requires root; device names are read from the running system.}]
daxctl list                      # one device per HBM NUMA node
daxctl reconfigure-device --mode=system-ram all
numactl --hardware               # HBM nodes now listed, without CPUs
\end{lstlisting}

On Lenovo's dual-socket example this involves two devices with Quadrant
clustering and eight with SNC4.  The conversion does not persist across
reboots, so it belongs in the node's administrator-managed start-up
configuration rather than in a user's job~\cite{lp1738}.
